\section{Methodology}
\label{sec:methodology}

This section sets out the baseline environment, the three friction extensions, and the unified factorial implementation. It also flags, in each subsection, every place where the implemented design departs from what the proposal specified, with a brief justification.\todo{Make sure the methodology is read in the order: baseline $\to$ each friction $\to$ unified design. The deviations from the proposal --- noise distribution, async schedule, off-path tests --- are flagged inline; consider whether to also collect them in a single ``deviations from the proposal'' table at the end of this section.}

\subsection{Baseline environment}
\label{sec:methodology-baseline}

The baseline reproduces \citet{calvano2020}. Two firms (\(n=2\)) repeatedly play a Bertrand pricing game with logit demand. The action set is a discrete grid of \(m=15\) prices, constructed by taking the static-Bertrand Nash and the (symmetric) joint-monopoly prices and extending the range slightly above and below them. Demand has the standard logit form
\begin{equation}
    q_i(p_1, p_2) \;=\; \frac{\exp\!\left(\frac{a_i - p_i}{\mu}\right)}{\sum_{j=1}^{n}\exp\!\left(\frac{a_j - p_j}{\mu}\right) + \exp\!\left(\frac{a_0}{\mu}\right)},
    \label{eq:logit-demand}
\end{equation}
with vertical-quality parameters \(a_1 = a_2 = 2\), outside-option \(a_0=0\), and horizontal-differentiation parameter \(\mu = 0.25\). Marginal costs are \(c_1 = c_2 = 1\). Single-period profits are \(\pi_i(p_1,p_2) = (p_i - c_i)\,q_i(p_1,p_2)\).

The state at time \(t\) is the vector of last-period prices \(s_t = (p_{1,t-1}, p_{2,t-1})\), so memory length \(k=1\). With \(m\) price levels per agent the state space has \(m^{2}\) elements. Each agent maintains a tabular action-value function \(Q_i: \mathcal{S}\times\mathcal{A}_i \to \R\). Actions are selected by an \(\varepsilon\)-greedy rule with exponentially decaying exploration:
\begin{equation}
    a_{i,t} \;=\; \begin{cases} \mathrm{argmax}_{a}\, Q_{i,t}(s_t,a) & \text{with probability } 1-\varepsilon_t,\\ \mathrm{Uniform}(\mathcal{A}_i) & \text{with probability } \varepsilon_t,\end{cases} \qquad \varepsilon_t = e^{-\beta\,t}.
\end{equation}
The Q-table updates are the standard one-step Q-learning rule
\begin{equation}
    Q_{i,t+1}(s_t, a_{i,t}) \;=\; (1-\alpha)\,Q_{i,t}(s_t,a_{i,t}) \;+\; \alpha\!\left[\pi_i(a_t) + \delta\,\max_{a'}Q_{i,t}(s_{t+1}, a')\right],
    \label{eq:Q-update}
\end{equation}
with learning rate \(\alpha = 0.15\) and discount factor \(\delta = 0.95\). The exploration parameter is \(\beta = 4\times 10^{-6}\), implemented as a per-episode multiplier \(\varepsilon_{\text{decay}}=0.02\) over episodes of length \(5{,}000\) iterations.

A run is a sequence of one or more \emph{sessions}. Each session draws fresh random initial Q-values and initial prices, and is run until the greedy strategy is unchanged for four consecutive episodes (the standard convergence criterion in this literature) or a 1{,}000-episode cap is reached. Reported quantities are means across sessions of session-level converged values. The primary outcome variable is the average profit gain
\begin{equation}
    \Delta \;=\; \frac{\bar{\pi} - \Pnash}{\Pmono - \Pnash},
\end{equation}
which equals \(0\) when the agents play the static-Bertrand Nash equilibrium and \(1\) when they play the symmetric joint-monopoly prices. We also report the cross-session standard deviation of session-level \(\Delta\) (denoted SD throughout) as a measure of dispersion.

\paragraph{Implementation and validation.}
The simulator is a from-scratch C++17 translation of the Calvano Fortran reference code. The translation was validated against the Fortran output at \(m=15\) before any frictions were introduced, and produces \(\Delta = 0.848\) (1{,}000 sessions, all converged), within sampling noise of \citet{calvano2020}'s reported value of approximately \(0.84\). A grid-robustness check at \(m\in\{15, 25, 50, 100\}\) confirms the result is not a coarse-grid artefact: \(\Delta\) declines mildly from $0.848$ at \(m=15\) to \(0.706\) at \(m=100\) and is essentially flat between \(m=50\) and \(m=100\). All runs use Calvano's original fixed-\(\beta\) protocol; the \(\nu\)-corrected version that holds learning quality constant across \(m\) is left for future work.

\subsection{Friction 1: observation latency}
\label{sec:methodology-latency}

Agent \(i\) at time \(t\) observes the rival's price from period \(t-1-L\) instead of \(t-1\). The own-price observation is unchanged. The agent's perceived state is therefore \(\tilde{s}^{(i)}_t = (p_{i,t-1},\, p_{j,t-1-L})\), which differs across agents whenever \(L\geq 1\). All other elements of the algorithm --- action selection, reward, Q-update, convergence test --- are applied to the per-agent perceived state. Profits are computed from the \emph{true} joint action, so the friction degrades information without altering payoffs.

A per-agent FIFO buffer of length \(L+1\) tracks each agent's price history, initialised by repeating each agent's randomly drawn initial price across the buffer slots. Because exploration is near-uniform during the first \(L\) iterations of every session (\(\varepsilon\approx 1\)) the buffer initialisation has negligible effect. At \(L=0\) the buffer reduces to a single slot and the entire learner reproduces the baseline bit-for-bit, providing an internal sanity check that confirms correctness of the implementation before each parameter sweep.

\subsection{Friction 2: noisy monitoring}
\label{sec:methodology-noise}

Each period, each agent observes the rival's previous-period price corrupted by an independent Gaussian draw on the \emph{price index}:
\begin{equation}
    \tilde{p}^{(i)}_{j,t-1} \;=\; \mathrm{clamp}\!\left(\,\mathrm{round}\!\left(k_{j,t-1} + \sigma\,\varepsilon^{(i)}_t\right),\; 1,\; m\right),
    \qquad \varepsilon^{(i)}_t \sim \mathcal{N}(0,1),
    \label{eq:noise}
\end{equation}
where \(k_{j,t-1}\in\{1,\dots,m\}\) is the rival's true price index and \(\sigma\) is the noise scale. The agent uses the noisy observation as the rival-side coordinate of its perceived state and cannot ``denoise''; noise is drawn independently per agent and per period. Profits are computed from the true joint action.

The proposal originally specified a discrete-flip noise model (probability \(q\) of perturbation to a neighbouring grid point); the implemented Gaussian-on-index version is a smooth one-parameter generalisation, locality-preserving, and reduces to the no-noise baseline at \(\sigma=0\). At \(\sigma=0\) the noise random number generator is short-circuited so the baseline is reproduced bit-for-bit, an internal sanity check parallel to the latency one.

\subsection{Friction 3: asynchronous actions}
\label{sec:methodology-async}

This friction implements the Maskin--Tirole alternating-move structure with a tunable lock-in length \(T\). At every iteration \(t\geq 1\) exactly one agent ``moves'' (revises its price) and the other holds its previous-period price. The mover identity follows a strict deterministic schedule:
\begin{equation}
    \mathrm{mover}(t) \;=\; \left(\!\left\lfloor \tfrac{t-1}{T} \right\rfloor \!\bmod\! 2\right) + 1.
\end{equation}
At \(T=0\) the rule is short-circuited and both agents move every period, recovering the simultaneous baseline. At \(T=1\) the agents strictly alternate (the standard Maskin--Tirole case). At \(T=k\) the moving agent is locked in for \(k\) consecutive periods before the other agent takes over.

Both agents Q-learn every period regardless of who moved. The non-mover's ``action'' is its previous-period price (a valid element of its action set), and its Q-update uses the observed reward from the actual joint action just like the mover's. The exploration parameter \(\varepsilon_t\) decays for both agents every period, on the convention that one iteration is one unit of calendar time for everyone.

The proposal originally suggested random alternation plus off-path deviation tests; the implemented version uses strict deterministic alternation and varies \(T\) on a grid. Off-path deviation analysis is left for future work.

\subsection{Combined frictions}
\label{sec:methodology-combined}

To test whether the three frictions interact, a \emph{unified} learner accepts \((L, \sigma, T)\) jointly and applies all three simultaneously. Within each iteration the order of operations is:
\begin{enumerate}[itemsep=0pt, topsep=2pt]
    \item Each agent's perceived state is \(\tilde{s}^{(i)}_t = (p_{i,t-1},\, \mathrm{noise}(p_{j,t-1-L}))\) --- noise is applied to the already-lagged observation.
    \item The mover (under \(T\)-period lock-in) is determined; the non-mover's action equals its previous-period price.
    \item The mover selects an action via \(\varepsilon\)-greedy on its own perceived state.
    \item Profits are computed from the true joint action.
    \item Both agents Q-update on their own perceived current and next-period state.
    \item Price-history buffers and \(\varepsilon\) decay one step.
\end{enumerate}
At \((L=0,\sigma=0,T=0)\) the unified learner reproduces the baseline. At single-friction corners it reproduces the corresponding single-friction learner to numerical precision. These two correspondences serve as built-in cross-validation tests of the unified implementation against each of the three single-friction implementations.

\paragraph{Factorial design.}
The factorial uses two levels per friction: ``off'' (the friction's null value) and ``on'' (a moderate, non-saturating value). The on-levels are chosen so that each single-friction effect is visible but not at the floor:
\begin{itemize}[itemsep=0pt, topsep=2pt]
    \item \(L \in \{0, 2\}\) --- at \(L=2\) the latency-only \(\Delta\) is approximately \(0.215\) (a \(75\%\) reduction).
    \item \(\sigma \in \{0, 1\}\) --- at \(\sigma=1\) the noise-only \(\Delta\) is approximately \(0.493\) (a \(42\%\) reduction).
    \item \(T \in \{0, 1\}\) --- at \(T=1\) the async-only \(\Delta\) is approximately \(0.378\) (a \(55\%\) reduction).
\end{itemize}
The eight cells are labelled \texttt{F$abc$} where each digit \(a,b,c\in\{0,1\}\) indicates whether \(L\), \(\sigma\), \(T\) is on. Each cell is run at \(N=250, 500, 1000\) sessions to verify that the interaction pattern is not an artefact of replication noise.

\paragraph{Session counts and standard errors.}
The single-friction sweeps and the \(N=250\) factorial use 250 sessions per cell. The factorial replications add 500 and 1000 sessions per cell. The reported SD is the standard deviation of session-level \(\Delta\) values across sessions; the standard error of the mean scales as \(\mathrm{SD}/\sqrt{N}\), and is below \(0.02\) on the headline factorial cells at \(N=1000\), tight enough to distinguish the three benchmark predictions in Section~\ref{sec:results-factorial}.
